\section*{Bab \ref{ch:b2:meiosis-heredity} --- Meiosis, Pencampuran Genetik, dan Pewarisan}

\begin{solution}{exo:b2:meiosis-heredity:1}
Mitosis: satu pembelahan, tanpa pemasangan homolog, dua sel, yang
masing-masing serupa secara genetik dengan induknya ($2n \to 2n$).
\omterm{def:b2:meiosis-heredity:meiosis}{Meiosis}: dua pembelahan sesudah satu replikasi, homolognya berpasangan
dan berpindah silang, empat sel, \omterm{def:b2:meiosis-heredity:meiosis}{haploid} dan semuanya berbeda secara
genetik ($2n \to n$).
\end{solution}

\begin{solution}{exo:b2:meiosis-heredity:2}
$2^{23} = 8.4\times 10^{6}$; $2^{4} = 16$; $2^{7} = 128$ --- sebelum
\omterm{prop:b2:meiosis-heredity:stages}{pindah silang} melipatgandakan tiap angka itu tanpa batas.
\end{solution}

\begin{solution}{exo:b2:meiosis-heredity:3}
$Tt\times Tt$: \omterm{def:b2:meiosis-heredity:vocabulary}{genotipe} $1\,TT : 2\,Tt : 1\,tt$, \omterm{def:b2:meiosis-heredity:vocabulary}{fenotipe} $3$ tinggi
: $1$ kerdil. $Tt\times tt$: $1\,Tt : 1\,tt$, yaitu $1 : 1$.
Silangujikan tumbuhan yang tinggi itu kepada tumbuhan kerdil: semuanya
tinggi berarti $TT$, separuhnya kerdil berarti $Tt$ (atau serbuki
sendiri: $TT$ berbiak sejati, $Tt$ memilah $3 : 1$).
\end{solution}

\begin{solution}{exo:b2:meiosis-heredity:4}
Gen yang terpaut terletak pada kromosom yang sama dan gabungan \omterm{def:b2:meiosis-heredity:vocabulary}{alel}
tetuanya berlebihan di antara gametnya; \omterm{def:b2:meiosis-heredity:linkage}{frekuensi rekombinasinya}
adalah bagian gamet rekombinan (keturunan rekombinan sebuah persilangan
uji); satu \omterm{def:b2:meiosis-heredity:linkage}{sentimorgan} adalah rekombinasi satu persen.
Gen yang berjauhan pada kromosom yang panjang memiliki sedikitnya satu
\omterm{prop:b2:meiosis-heredity:stages}{pindah silang} di antaranya pada hampir tiap \omterm{def:b2:meiosis-heredity:meiosis}{meiosis}, dan \omterm{prop:b2:meiosis-heredity:stages}{pindah silang}
berganda mengacak gabungannya, sehingga $r$ mencapai $\tfrac12$ ---
tak terbedakan dari \omterm{thm:b2:meiosis-heredity:mixing}{pemilahan bebas}.
\end{solution}

\begin{solution}{exo:b2:meiosis-heredity:5}
Jumlahnya 556; harapannya $312.75 : 104.25 : 104.25 : 34.75$; $\chi^2 =
2.25^2/312.75 + 3.75^2/104.25 + 3.25^2/104.25 + 2.75^2/34.75 = 0.016
+ 0.135 + 0.101 + 0.218 = 0.47 < 7.81$: sepadan dengan $9 : 3 : 3 :
1$.
\end{solution}

\begin{solution}{exo:b2:meiosis-heredity:6}
Terhadap $1 : 1 : 1 : 1$ (250 tiap golongan): $\chi^2 = (162^2 + 138^2 + 147^2
+ 153^2)/250 = 361 \gg 7.81$: terpaut. Rekombinannya $103 + 97 = 200$
dari 1000: $r = 0.20$, yaitu \qty{20}{cM}. Dalam keadaan tolak, $Ab/aB$
memberi 400 $Ab$, 400 $aB$, 100 $AB$, 100 $ab$: jaraknya sama, hanya
golongan tetua dan rekombinannya yang bertukar.
\end{solution}

\begin{solution}{exo:b2:meiosis-heredity:7}
$d = -\tfrac12\ln(1 - 2r)$: $r = 0.20$ memberi $-\tfrac12\ln 0.6 =
0.255$, yaitu \qty{25.5}{cM}; $r = 0.45$ memberi
$-\tfrac12\ln 0.1 = 1.15$, yaitu \qty{115}{cM}.
$r = \tfrac12(1 - e^{-2d})$: \qty{30}{cM} memberi
$\tfrac12(1 - e^{-0.6}) = 0.226$; \qty{100}{cM} memberi $\tfrac12(1 -
e^{-2}) = 0.43$. Di bawah \qty{10}{cM} pembetulannya
dapat diabaikan; di atas \qty{30}{cM} ia penting, dan di atas
\qty{50}{cM} $r$ menjenuh sehingga sebuah persilangan dua titik tunggal
tidak lagi dapat mengukur jaraknya --- gen yang berjauhan dipetakan
dengan merangkai gen yang berdekatan.
\end{solution}

\begin{solution}{exo:b2:meiosis-heredity:8}
(a) $X^wX^w \times X^+Y$: semua anak jantannya putih ($X^wY$), semua
anak betinanya merah ($X^+X^w$). (b) $X^+X^w \times X^wY$: anak
jantannya separuh merah, separuh putih; anak betinanya separuh merah
($X^+X^w$), separuh putih ($X^wX^w$). (c) anak betina $X^+X^w$
$\times$ saudara jantannya $X^wY$: seperti pada (b) --- separuh tiap
kelaminnya putih, separuhnya merah.
\end{solution}

\begin{solution}{exo:b2:meiosis-heredity:9}
Dua enzim yang berderet membuat pigmennya; ungu membutuhkan sebuah \omterm{def:b2:meiosis-heredity:vocabulary}{alel}
\omterm{def:b2:meiosis-heredity:vocabulary}{dominan} pada kedua lokusnya ($C\_P\_$). Galurnya adalah $CCpp$ dan
$ccPP$ (masing-masing putih dengan alasan yang berbeda); hibridnya
$CcPp$ berwarna ungu; kalau diserbuki sendiri ia memberi $9\,C\_P\_$
ungu berbanding $7$ putih ($3\,C\_pp + 3\,ccP\_ +
1\,ccpp$). Persilangan uji $CcPp\times ccpp$: $1$ ungu ($CcPp$) : $3$
putih.
\end{solution}

\begin{solution}{exo:b2:meiosis-heredity:10}
Tetuanya $+++$, $abc$; \omterm{prop:b2:meiosis-heredity:stages}{pindah silang} gandanya $+b+$, $a+c$; $b$ adalah
gen yang bertukar: urutannya $a$--$b$--$c$. $r_{ab} = (11 + 9 + 2)/1202
= 0.018$ (\qty{1.8}{cM}); $r_{bc} = (118 + 112 + 2)/1202 = 0.193$
(\qty{19.3}{cM}). Harapan gandanya $0.018\times 0.193\times 1202 =
4.2$; amatannya 2: kebetulannya $0.47$, interferensinya $0.53$.
\end{solution}

\begin{solution}{exo:b2:meiosis-heredity:11}
Askus pembelahan keduanya $300/1000 = \qty{30}{\%}$; jaraknya $=
\tfrac12\times 30 = \qty{15}{cM}$. Sebuah \omterm{prop:b2:meiosis-heredity:stages}{pindah silang} antara gennya
dan sentromernya hanya melibatkan dua dari empat kromatidnya, sehingga
askus yang memperlihatkan pemilahan pembelahan kedua mengusung dua
kromatid rekombinan dan dua kromatid tetua: \omterm{def:b2:meiosis-heredity:linkage}{frekuensi rekombinasinya}
adalah separuh frekuensi askus semacam itu. Bagi $2 : 4 : 2$: sebuah
\omterm{prop:b2:meiosis-heredity:stages}{pindah silang} antara lokusnya dan sentromernya mempertukarkan kedua
kromatid dalamnya, sehingga tiap setengah gelendongnya pada \omterm{def:b2:meiosis-heredity:meiosis}{meiosis} I
mengusung satu kromatid $A$ dan satu kromatid $a$; \omterm{def:b2:meiosis-heredity:meiosis}{meiosis} II lalu
memisahkannya dalam urutan $A\,a \mid
a\,A$ (atau sebaliknya), dan mitosis akhirnya melipatduakan tiap
sporanya: $2 : 4 : 2$.
\end{solution}

\begin{solution}{exo:b2:meiosis-heredity:12}
Ibunya adalah pembawa wajib (ia menerima satu-satunya X ayahnya),
sehingga perempuan itu menerima \omterm{def:b2:meiosis-heredity:vocabulary}{alelnya} dengan peluang $\tfrac12$.
Dua anak lelaki yang sehat: seorang pembawa memiliki peluang
$(\tfrac12)^2 = \tfrac14$ memperoleh dua anak lelaki yang sehat,
seorang bukan pembawa berpeluang 1. Bayes:
$P(\text{pembawa}\mid\text{data}) = (\tfrac12\times\tfrac14)/(\tfrac12
\times\tfrac14 + \tfrac12\times 1) = \tfrac{1/8}{5/8} = \tfrac15$.
\end{solution}

\begin{solution}{pb:b2:meiosis-heredity:1}
\textbf{1.} Tetuanya $vg/vg\;+/+$ dan $+/+\;e/e$; $F_1$-nya $+/vg\;+/e$;
gametnya $+\,+$, $+\,e$, $vg\,+$, $vg\,e$, masing-masing $\tfrac14$.
\textbf{2.} $900 : 300 : 300 : 100$.
\textbf{3.} $\chi^2 = 8^2/900 + 9^2/300 + 4^2/300 + 3^2/100 = 0.07 +
0.27 + 0.05 + 0.09 = 0.48 < 7.81$: \omterm{thm:b2:meiosis-heredity:mixing}{pemilahan bebasnya} berlaku.
\textbf{4.} Yang \omterm{def:b2:meiosis-heredity:vocabulary}{homozigot} liar pada kedua lokusnya adalah
$\tfrac{1}{16}$ dari semuanya, sedangkan $\tfrac{9}{16}$ bertipe liar:
jadi $\tfrac19$.
\textbf{5.} Liar, vestigial, eboni, vestigial eboni, masing-masing
$\tfrac14$.
\textbf{6.} \omterm{def:b2:meiosis-heredity:vocabulary}{Fenotipenya} adalah hasil kali \omterm{def:b2:meiosis-heredity:vocabulary}{genotipe} dan lingkungannya
(suhunya bekerja pada enzim pigmennya): sebuah \omterm{def:b2:meiosis-heredity:vocabulary}{genotipe} menetapkan
sebuah norma reaksi, bukan sebuah ciri yang tetap.
\textbf{7.} $b\,+/+\,vg$ (tolak): $b$ datang bersama $+_{vg}$ dari satu
tetua, $vg$ bersama $+_b$ dari tetua yang lain.
\textbf{8.} Tetuanya: hitam ($b\,+$) 965 dan vestigial ($+\,vg$) 944;
rekombinannya: hitam vestigial ($b\,vg$) 206 dan liar ($+\,+$) 185.
Kromosom tipe liar $+\,+$ tidak ada di dalam $F_1$-nya; ia hanya dapat
dibuat lewat sebuah \omterm{prop:b2:meiosis-heredity:stages}{pindah silang}.
\textbf{9.} $r = 391/2300 = 0.17$: \qty{17}{cM}.
\textbf{10.} $d = -\tfrac12\ln(1 - 0.34) = -\tfrac12\ln 0.66 = 0.21$:
\qty{21}{cM}.
\textbf{11.} Tidak ada \omterm{prop:b2:meiosis-heredity:stages}{pindah silang} pada \emph{Drosophila} jantan:
gamet jantannya hanya mengusung kedua kromosom tetuanya.
\textbf{12.} $F_1$ $b\,vg/+\,+$ (gandeng): golongan tetuanya liar dan
hitam vestigial, masing-masing \qty{41.5}{\%}; rekombinannya hitam dan
vestigial, masing-masing \qty{8.5}{\%}.
\textbf{13.} $X^hX^+$: seorang anak perempuan menerima satu-satunya X
ayahnya, yang mengusung $h$, dan sebuah X normal dari ibunya yang tidak
terkena.
\textbf{14.} $\tfrac12$ (anak lelaki) $\times$ $\tfrac12$ (menerima $X^h$) $=
\tfrac14$.
\textbf{15.} $\tfrac12$: ia menerima salah satu kromosom X ibunya;
ayahnya memberi sebuah X normal.
\textbf{16.} Tiap anaknya tidak terkena dengan peluang $\tfrac34$:
$(\tfrac34)^3 = \tfrac{27}{64} = 0.42$.
\textbf{17.} Seorang anak lelaki menerima Y ayahnya, tidak pernah
X-nya. Seorang perempuan terkena hanya sebagai $X^hX^h$: dari ayah
yang terkena dan ibu yang pembawa (atau terkena), yang jarang terjadi.
\textbf{18.} Peluang awalnya $\tfrac12$. Dua anak lelaki yang sehat
berpeluang $\tfrac14$ kalau ia pembawa, dan 1 kalau bukan: peluang
akhirnya $(\tfrac12\times
\tfrac14)/(\tfrac12\times\tfrac14 + \tfrac12) = \tfrac15$.
\textbf{19.} Satu-satunya X seorang anak lelaki datang dari ibunya, dan
Y-nya tidak mengusung satu pun gen ini, sehingga tiap anak lelakinya
menampakkan satu gamet ibunya tanpa tertutupi: persilangannya adalah
persilangan ujinya sendiri.
\textbf{20.} Tetuanya $+++$ (853), $ywm$ (833); \omterm{prop:b2:meiosis-heredity:stages}{pindah silang} gandanya
$+w+$ (2), $y+m$ (2); $w$ yang bertukar: urutannya $y$--$w$--$m$.
\textbf{21.} $r_{yw} = (24 + 26 + 2 + 2)/2000 = 0.027$; $r_{wm} = (128
+ 132 + 2 + 2)/2000 = 0.132$. Petanya: $y$ --- \qty{2.7}{cM} --- $w$ ---
\qty{13.2}{cM} --- $m$.
\textbf{22.} Harapannya $0.027\times 0.132\times 2000 = 7.1$; amatannya
4: kebetulannya $0.56$, interferensinya $0.44$.
\textbf{23.} Rekombinan $y$--$m$: $y++$, $+wm$, $yw+$, $++m$: $(24 +
26 + 128 + 132)/2000 = 0.155$, terhadap $0.027 + 0.132 = 0.159$: keempat
\omterm{prop:b2:meiosis-heredity:stages}{pindah silang} gandanya memulihkan gabungan tetua $y$--$m$ sehingga
terlewat, yaitu $2\times 4/2000 = 0.004$.
\textbf{24.} Anak betinanya menerima X $y\,w\,m$ ayahnya dan satu X
dari ibunya: semuanya \omterm{def:b2:meiosis-heredity:vocabulary}{heterozigot} dan bertipe liar secara \omterm{def:b2:meiosis-heredity:vocabulary}{fenotipe},
X rekombinan mana pun yang diusungnya; rekombinasinya tak terlihat.
\textbf{25.} $y$ --- \qty{2.7}{cM} --- $w$ --- \qty{13.2}{cM} --- $m$;
koefisien kebetulannya $0.56$, interferensinya $0.44$.
\end{solution}
